\studySubsection{practice-calc-06}{练习库：高等数学第 6 讲 辅助函数构造}

\problemAnchor{MATH1-CALC-0015}
\begin{problemBox}
\begin{problemMeta}
\textbf{题目编号：} \problemRef{MATH1-CALC-0015} \quad
\textbf{答案：} \answerRef{MATH1-CALC-0015} \quad
\textbf{来源：} Codex 原创 / K064 深度讲义配套练习（非真题）
\end{problemMeta}

\begin{enumerate}
  \item[\textbf{A}] 设 \(f\in C[0,1]\cap D(0,1)\)，且
  \(f(0)=f(1)=0\)。证明：存在 \(\xi\in(0,1)\)，使
  \[
  f'(\xi)+2f(\xi)=0.
  \]

  \item[\textbf{B}] 设 \(f\in C[1,2]\cap D(1,2)\)，且
  \(f(2)=4f(1)\)。证明：存在 \(\xi\in(1,2)\)，使
  \[
  \xi f'(\xi)=2f(\xi).
  \]

  \item[\textbf{C}] 设 \(f\in C[0,1]\)，记
  \[
  I=\int_0^1 2x f(x)\,\mathrm dx.
  \]
  证明：存在 \(\xi\in(0,1)\)，使 \(f(\xi)=I\)。

  \item[\textbf{M}] 设 \(\lambda\in\mathbb R\)，
  \(f\in C[0,3]\)，且 \(f\) 在 \((0,3)\) 内二阶可导。已知
  \[
  f(0)=f(3)=0,\qquad f(1)=\lambda.
  \]
  证明：存在 \(\xi\in(0,3)\)，使
  \[
  f''(\xi)=-\lambda.
  \]
\end{enumerate}
\end{problemBox}
